\section{Introduction}\label{sec:intro}

Program slicing finds the statements of a program that can affect a value at a point of interest, and
debugging is where developers use it~\cite{weiser1982programmers}. Quantum programs are hard to debug.
Their intermediate states cannot be inspected without disturbing them, their outcomes are probabilistic,
and the effect of a gate on one qubit may show only in the outcomes of
another~\cite{ali2022sequantum, murillo2025roadmap}. Developers who are new to quantum computing also find
the code itself hard to understand~\cite{daloisio2024llm}. A slicer that tells a developer which lines of a
quantum program can influence a wrong outcome would help with both problems.

Slicing has been proposed for quantum programs. Yu and Zhao's quantum program dependence
graph~\cite{yu2025qpdg} extends the classical program dependence graph~\cite{ferrante1987pdg} with
entanglement, no-cloning, temporal and measurement dependences, and reduces slicing to graph reachability
under the classical criterion $(s, V)$. A roadmap of quantum software engineering presents this graph as
the basis for slicing, debugging and testing~\cite{murillo2025roadmap}, and a later one notes that
validating such models remains open~\cite{ye2026llmqse}. No definition of quantum dependence has been
checked against the semantics of quantum programs. It is therefore not known whether the slices such
definitions produce contain the operations a developer must inspect.

Section~\ref{sec:example} shows a real bug where this matters. Reading a two-qubit gate as a classical
assignment gives a slice that omits the faulty line and still reproduces the qubit's state, and the bug
itself shows only in the correlation between two qubits, which a single-variable criterion cannot
express.

This paper defines quantum dependence semantically, as the set of operations where a change can alter
the state of the criterion qubits for some input, and generalizes the criterion to a set of qubits whose
joint state the slice must preserve. We then measure how dependence definitions behave against this
ground truth, in a pre-registered study on \numNPairs{} circuit--qubit pairs from three benchmark suites,
\numFamCircuits{} circuits with known structure, \numRqThreeUnits{} seeded faults, and real bugs from
Bugs4Q~\cite{zhao2023bugs4q}. The definitions range from the classical program dependence graph to
pruning by simulation, and they include our documented reading of Yu and Zhao's graph.

A correct slice needs data, control and quantum dependence. The classical program dependence graph
has the first two, and it still gives wrong slices for \numCpdgUnsoundPct{} of the pairs. Control
dependence repairs it on circuits with measurements and feedback, where it is wrong on
\numCpdgDynamicUnsound{} of \numNDynamicPairs{} pairs against \numDuDynamicUnsound{} without control
dependence. It cannot repair the rest, because a two-qubit gate can change its control qubit (phase
kickback), a quantum dependence that no classical reading of the gate expresses. Following entanglement, as the dependence
graph does, is sound but keeps \numRqOneYzSize{} of the circuit. A simple rule, in which every gate reads
and writes all its qubits, keeps \numRqOneLcSize{}, only \numLcExcess{} points above the dependences
themselves, misses no dependence, and stays correct on every
tested input, and its soundness is machine-checked. Entanglement turns out to be neither necessary nor
sufficient for dependence. Pruning by commutation shrinks the slices further and stays correct on all
random inputs. Pruning by simulation shrinks them by \numHTwoRelative{}, but only for the input it was
computed on. It removes the faulty operation in \numCQFour{} seeded faults and lowers the
fault-localization accuracy of three language models by up to \numTTwoMaxDrop{} percentage points. A
criterion over all affected qubits recovers \numCQFourRecovered{} of these faults.

This paper makes four contributions. It gives semantic definitions of quantum dependence, slice
soundness and dependence completeness, with a criterion over sets of qubits, and proves that the wire
rule is sound and complete (Section~\ref{sec:dependence}). It reports the first comparison of quantum
dependence definitions against a semantic ground truth, covering soundness, completeness, robustness to
inputs and precision. It measures what each definition does to fault localization, by hand, with
mutation-based fault localization and with language models, on seeded faults and on real bugs mapped to
the lines of their Qiskit programs. Finally, it provides a replication package that regenerates every
number in the paper.

Section~\ref{sec:example} presents the debugging scenario, Section~\ref{sec:dependence} the definitions,
Section~\ref{sec:design} the study design and Section~\ref{sec:results} the results.
Section~\ref{sec:discussion} discusses what the results mean for tool builders,
Section~\ref{sec:threats} the threats to validity, Section~\ref{sec:related} related work, and
Section~\ref{sec:conclusion} concludes.
